\section{Notations and backgrounds}
\subsection{Symmetric groups}
Recall that the symmetric group is the Weyl group of the special linear group. Define the usual action of the symmetric group $S_m$ on $\mathbb{N}^m$ by $ \sigma(x_1,\dots,x_m)= (x_{\sigma(1)}, \dots, x_{\sigma(m)})$. For any $A\in \mathbb{N}^m$, define $ H_A \le S_m $ to be the subgroup $ \{\sigma \in S_m\colon \sigma(A)=A\}$. Let $ D_A \subset S_m $ be the set of (left) coset representatives of $H_A$ with the fewest inversions: in other words, $ \sigma \in D_A $ if and only if $ \inv (\sigma) \le \inv (\tau) $ for every $ \tau \in \sigma H_A$. Here, $\inv$ is the length function defined as the fewest number
of simple reflections constituting the Weyl group element.

 \begin{Example} Take $ A= (3,3,3,2,2,1).$ Then $H_A=S_3 \times S_2 \times S_1 \leq S_6$. We have $ (34) \in D_A $ with $ (34) \cdot A= (3,3,2,3,2,1)$, and $ (35) \in D_A $ with $ (35)\cdot A=(3,3,2,2,3,1) $. However, $ (134)\cdot A= (34) \cdot A$ and $\inv((134)) > \inv((34))$, so $ (134) \notin D_A $.
 \end{Example}

 We recall (see, e.g., \cite{Carter}) that each coset of $H_A$ has a unique representative $\sigma \in D_A$, and that $\inv(\sigma\tau) = \inv(\sigma) + \inv(\tau)$ for every $\tau \in H_A$.

 \begin{Lemma}\label{Gene}
 	Suppose that $\tau,\tau' \in D_A$. Then there exist a sequence of elements $\tau_0,\dots,\tau_l\in D_A$ such that $\tau_0 = \tau'$, $\tau_l=\tau$ and every $\tau_{j+1}\tau_j^{-1}$ is a transposition for $j=0,\dots,l-1$.
 \end{Lemma}
 \begin{proof}
 	It suffices to prove this statement when either $\tau$ or $\tau'$ is the identity permutation $e$, because the two sequences can be concatenated. So suppose that $\tau$ is an arbitrary element of $D_A$ and $\tau'=e$. Let $s_k\dots s_1$ be a minimal word representation of $\tau$ and set $\tau_j = s_j \dots s_1$. Then $\tau_{j+1}\tau_j^{-1}=s_{j+1}$, which is a transposition. So it remains to show that $\tau_j \in D_A$. If it were not, then there would exist a transposition $s\in H_A$ such that $\inv(\tau_js) = \inv(\tau_j)-1$. But then $\inv(\tau s) = \inv(\tau)-1$, contradicting the assumption that $\tau \in D_A$.
 	
 	Now suppose that $\tau=e$ and $\tau'$ is an arbitrary element of $D_A$. By the previous paragraph, there exist a sequence of elements $\tilde{\tau}_0=e,\tilde{\tau}_1,\dots,\tilde{\tau}_l=\tau'$ in $D_A$ such that every $\tilde{\tau}_{j+1}\tilde{\tau}_j^{-1}$ is a~transposition. Setting $\tau_j = \tilde{\tau}_{l-j}$, we have that $\tau_0=\tau',\dots,\tau_l=e$ is a sequence of elements in $D_A$ and ${\tau}_{l-j-1}{\tau}_{l-j}^{-1}$ is a transposition for every $j$. The latter equality is equivalent to the condition that for every $k$, $\tau_k = s\tau_{k+1}$ for some transposition $s$. Since this is also equivalent to the condition that $\tau_{k+1}\tau_k^{-1}$ is a transposition for every $k$, this finishes the proof.
 \end{proof}

 Let $ \mathcal{B}_m^{(N)} $ denote the set of sequences $ \mu=(\mu_0,\dots,\mu_N) $ of non-negative integers such that $ \mu_0+\dots+\mu_N=m $. For any $ \mu \in \mathcal{B}_m^{(N)} $, let $ H^\mu \le S_m$ denote the subgroup $ S_{\mu_0}\times \dots \times S_{\mu_N} $, and likewise let $ D^\mu $ denote the set of left coset representatives with the fewest inversions. Define $\mathcal{B}_m$ to be the union $\bigcup_{N\geq 1} \mathcal{B}_m^{(N)}$.

Let $\mathcal{W}_m \subset \mathbb{N}^m$ denote the subset of elements $(i_1,\dots,i_m)$ satisfying $i_1 \leq \dots \leq i_m$. For $\ii \in \mathcal{W}_m$, and assuming that $N\geq i_m$, let $\mu^{(N)}(\ii) \in \mathcal{B}_m^{(N)}$ be defined by
\[
\mu^{(N)}(\ii)_k = \vert \{ l \in \{1,\dots,m\}\colon i_l=k \} \vert.
\]
For $\ii \in \mathcal{W}_m$ satisfying $i_m \leq N$, we have a natural isomorphism between the subgroups $H_{\ii}$ and $H^{\mu^{(N)}(\ii)}$. Thus there is also a natural bijection between $D_{\ii}$ and $D^{\mu^{(N)}(\ii)}$.

Given $\ii \in \mathcal{W}_m$, define the equivalence relation $\sim_{\ii}$ on $S_m$ by
\[
\tau \sim_{\ii} \sigma \text{ if and only if } \sigma^{-1}\tau \in H_{\ii}.
\]
In words, $\tau \sim_{\ii} \sigma$ if and only if they are in the same left coset of $H_{\ii}$. For any $\tau \in S_m$, there is a unique $\sigma \in D_{\ii}$ such that $\tau \sim_{\ii} \sigma$. Define $d_{\ii}(\tau)= \inv\big(\sigma^{-1}\tau\big)$. In other words, if $\tau$ is written uniquely as $\tau=\sigma \xi$ for $\sigma \in D_{\ii}$ and $\xi \in H_{\ii}$, then $d_{\ii}(\tau) = \inv(\xi)$. We will also let $\sigma_{\ii}(\tau)$ and $\xi_{\ii}(\tau)$ denote the two permutations in the unique expression $\tau = \sigma \xi$.

Finally, given $\tau \in S_m$, let $\bar{\tau}$ denote the reversed permutation $\bar{\tau}(k)= \tau(m+1-k)$.

We conclude this section by noting the following identity.
\begin{Lemma}\label{Leem}
For $\ii \in \mathcal{W}_m$, set $\mu = \mu^{(N)}(\ii)$. Then
\[
-N\mu_{0} -(N-2)\mu_{1} + \dots + N\mu_{N} = -Nm + 2(i_1+ \dots + i_m).
\]
\end{Lemma}
\begin{proof}
By definition,
\[
\mu_k = | l \in \{1,\dots,m\}\colon i_l=k |.
\]
We thus rewrite
\begin{align*}
-N\mu_{0} -(N-2)\mu_{1} + \dots + N\mu_{N} &= -N(\mu_0 + \dots + \mu_N) + 2\mu_1 + 4\mu_2 + \dots + 2N\mu_N \\
&= -Nm+ 2\mu_1 + 4\mu_2 + \dots + 2N\mu_N
\end{align*}
and note that
\[
2(i_1 + \dots + i_m) = 2\sum_{k=1}^N k\cdot | l \in \{1,\dots,m\}\colon i_l=k | = 2\sum_{k=1}^N k\cdot \mu_k,
\]
which shows the identity.
\end{proof}

\subsection{Quantum groups}
We use the following notation modified from \cite{Jimbo1986}.
Define $ U_q(\mathfrak{sl}(N+1)) $ to be the associative algebra over $ \mathbb{C} $ generated by the symbols $ q^{\pm h_i/2}$, $\hat{e}_{\pm, i}$, $1 \leq i\leq N$, under the following relations:
\begin{gather*}
q^{h_i/2} \cdot q^{-h_i/2}= q^{-h_i/2} \cdot q^{h_i/2}=1, \qquad q^{h_i/2} \cdot q^{h_{i'}/2}= q^{h_{i'}/2} \cdot q^{h_i/2},\nonumber\\
 q^{h_i/2} \hat{e}_{\pm, i'} q^{-h_i/2}= q^{\pm a_{ii'}/2} \hat{e}_{\pm, i'},\nonumber\\
[\hat{e}_{+,i}, \hat{e}_{-,i'}] =\delta_{i,i'} \frac{q^{h_i}-q^{-h_i}}{q-q^{-1}}, \nonumber\\
 \hat{e}_{\pm, i}\hat{e}_{\pm, i'} = \hat{e}_{\pm ,i'}\hat{e}_{\pm ,i}, \qquad |i-i'|\geq 2, \nonumber\\
 \hat{e}_{\pm, i}^2 \hat{e}_{\pm, i\pm 1} -\big(q+q^{-1}\big) \hat{e}_{\pm, i} \hat{e}_{\pm ,i\pm 1} \hat{e}_{\pm ,i} + \hat{e}_{\pm, i\pm 1} \hat{e}_{\pm, i}^2=0,\qquad 1\leq i, i\pm 1 \leq N. %\label{UqglDef}
\end{gather*}
Here, $ (a_{i,i'})_{1\leq i,i' \leq N} $ denotes the Cartan matrix of type $ A_N $, i.e.,
\begin{equation*}
a_{ii'}=
\begin{cases}
\hphantom{-}2,& i = i',\\ -1, & i=i' \pm 1,\\ \hphantom{-}0,& \text{otherwise}.
\end{cases}
\end{equation*}
Then define $ U_q(\mathfrak{gl}(N+1)) $ by adjoining to $ U_q(\mathfrak{sl}(N+1)) $ the elements $ q^{\pm \epsilon_i/2}$, $0 \leq i \leq N $, so that $ q^{h_i}= q^{\epsilon_{i-1}-\epsilon_i} $ and that $ q^{\epsilon_0+\dots +\epsilon_N} $ belongs to the center.

The $m$-fold co-product is the algebra homomorphism \[\Delta^{(m)}\colon \ \Uqgl \rightarrow \underbrace{\Uqgl \otimes \dots \otimes \Uqgl}_{m} \]
such that \begin{gather}
 \Delta^{(m)}\big (q^{\pm \epsilon_i/2}\big)= q^{\pm \epsilon_i/2} \otimes \dots \otimes q^{\pm \epsilon_i/2},\qquad
 \Delta^{(m)} (\hat{e}_{\pm ,i} ) =\sum_{v=1}^{m} \hat{e}^{(v)}_{\pm ,i} ,\label{qerel}
\end{gather}
where
\[
\hat{e}^{(v)}_{\pm ,i} = \underbrace{q^{h_i/2} \otimes \dots \otimes q^{h_i/2}}_{v-1} \otimes \hat{e}_{\pm ,i} \otimes \underbrace{ q^{-h_i/2} \otimes \dots \otimes q^{-h_i/2}}_{m-v}.
\]
 We also have the reversed co-product
 \begin{gather*}
 \bar{\Delta}^{(m)} \big(q^{\pm \epsilon_i/2}\big)= q^{\pm \epsilon_i/2} \otimes \dots \otimes q^{\pm \epsilon_i/2}, \\
 \bar{\Delta}^{(m)} (\hat{e}_{\pm, i})= \sum_{v=1}^{m} \underbrace{q^{-h_i/2} \otimes \dots \otimes q^{-h_i/2}}_{v-1} \otimes \hat{e}_{\pm, i} \otimes q^{h_i/2} \otimes \dots \otimes q^{h_i/2}.
 \end{gather*}
We write $\Delta$, $\bar{\Delta}$ for $\Delta^{(2)}$, $\bar{\Delta}^{(2)}$. The map $ \Delta $ endows $ \Uqgl $ with a structure of a bi-algebra. It is also a Hopf algebra, but we will not need the counit and antipode.

Consider the following relations for $R$:
\begin{gather}
 R \Delta(u) = \bar{\Delta}(u) R \qquad \forall u\in \Uqgl. \label{Equi}
\end{gather}
 This admits a unique (up to a multiplicative constant) solution
$ R \in \End (V\otimes V) $, where $ V:=\mathbb{C}^{N+1} $ is the defining representation of $ \Uqgl $.

 Let us consider \eqref{Equi} in $ \Uqgl \otimes \End \big(\C^{N+1}\big) $. Then a family of solutions is given by (see \cite{Jimbo1986}) $R(x)= \sum_{0 \leq i,j \leq N} \hat{E}_{ij}'(x) \otimes e_{ji} $, where
 \begin{gather*} %\label{ehatx}
 \hat{E}'_{ij}(x) = \begin{cases}
 \big(x q^{(\epsilon_i+\epsilon_j-1)/2}\big)^{\mp 1} {E}'_{ij} , & i \lessgtr j,\\
 \big(xq^{\epsilon_i}- x^{-1}q^{- \epsilon_i}\big)/\big(q-q^{-1}\big), & i=j.
 \end{cases}
 \end{gather*}
 For this paper, only the value at $x=1$ will be used in the proofs, but for completeness Jimbo's result is stated in its full generality. Here, $ E'_{ij} $ are the root vectors defined recursively by
 \begin{gather*}
E'_{ij}= E_{ik}'E_{kj}' -q^{\pm 1} E_{kj}'E_{ik}',\qquad i \lessgtr k \lessgtr j, \qquad E'_{i-1,i} = \hat{e}_{+,i}, \qquad E'_{i,i-1} = \hat{e}_{-,i}
\end{gather*}
and $e_{ji}$ is the usual matrix which acts on the canonical basis $\{I_0,\dots,I_N\}$ of $\mathbb{C}^{N+1}$ by
\[
e_{ji}I_k = \delta_{ik} I_j.
\]

In \cite{GZB}, the authors define
\[
\hat{E}_{ij}
=
\begin{cases}
q^{(E_{ii}+E_{jj}-1)/2}E_{ij}, & i \neq j,\\
q^{E_{ii}}/\big(q-q^{-1}\big), & i=j,
\end{cases}
\]
where the modified root vectors are
\begin{equation*}
E_{ij}= E_{ik}E_{kj} -q^{- 1} E_{kj}E_{ik},\qquad i \lessgtr k \lessgtr j, \qquad E'_{i-1,i} = \hat{e}_{+,i}, \qquad E'_{i,i-1} = \hat{e}_{-,i}.
\end{equation*}
Using Jimbo's results, they show that
\begin{gather*}
R = \sum_{N \geq i\geq j \geq 0} \hat{E}_{ij} \otimes e_{ji} \in \Uqgl \otimes \End (V),\nonumber\\
R^T = \sum_{N \geq i\geq j \geq 0} \hat{E}_{ji} \otimes e_{ij} \in \Uqgl \otimes \End (V) %\label{r}
\end{gather*}
satisfy
\begin{gather*}
R\Delta(u) = \bar{\Delta}(u) R,\qquad
R^T\bar{\Delta}(u) = {\Delta}(u) R^T.
\end{gather*}

We can solve \eqref{Equi} even more generally. Consider the $ m $-fold tensor of the defining representation $ V^{\otimes m}$, and let $P_mV^{\otimes m}$ be the symmetric projection. In other words, $P_mV^{\otimes m}$ is isomorphic to the module of highest weight~$m\nu$, where~$\nu$ is the highest weight of~$V$.

Then the solution of \eqref{Equi} in $ \Uqgl \otimes \End (P_m V^{\otimes m}) $ is given by the fused $R$-matrix~\cite{Jimbo1986}%
\begin{gather*} %\label{fusx}
R^{}_{0,\{1,2,\dots,m\}}(x)= R_{0m}(x)R_{0m-1}(xq) \cdots R_{01}\big(xq^{ m-1}\big) P_m \in \Uqgl\otimes \End \big(P_m V^{\otimes m}\big) .
\end{gather*}

Therefore, the fused $R$-matrices
\begin{gather*}
R^{}_{0,\{1,2,\dots,m\}} = R_{0m}R_{0m-1}\cdots R_{01}P_m \in U_q(\mathfrak{gl}_{N+1})\otimes \End \big(P_m V^{\otimes m}\big), \\
R^{T}_{0,\{1,2,\dots,m\}} = R^T_{0m}R^T_{0m-1}\cdots R^T_{01}P_m \in U_q(\mathfrak{gl}_{N+1})\otimes \End \big(P_m V^{\otimes m}\big)
\end{gather*}
satisfy
\begin{gather}
R_{0,\{1,2,\dots,m\}} \Delta(u) = \bar{\Delta}(u) R_{0,\{1,2,\dots,m\} },\qquad
R^T_{0,\{1,2,\dots,m\}} \bar{\Delta}(u) = {\Delta}(u) R^T_{0,\{1,2,\dots,m\}}\label{RRel}
\end{gather}
in $\Uqgl \otimes \End (P_m V^{\otimes m})$ for all $u\in \Uqgl$.

Therefore,
\begin{equation*}
\Gamma_m: = R^T_{0,\{1,2,\dots,m\}} R_{0,\{1,2,\dots,m\}} \in \Uqgl \otimes \End \big(P_m V^{\otimes m}\big)
\end{equation*}
commutes with $ \Delta (u) $ for all $u\in \Uqgl$. By Drinfeld's central element construction \cite{Drin90}, the element
\begin{equation*}
C_m= \id \otimes\tr_q(\Gamma_m)
\end{equation*}
is central in $\Uqgl$, where the quantum trace $ \tr_q $ of an operator $A$ is defined by
\begin{equation*}%\label{key}
\tr_q(A) = \tr\big(q^{-2h_{\rho}}A\big) := \tr \big(q^{-N\epsilon_0-(N-2)\epsilon_1 - \dots +N\epsilon_N}A\big).
\end{equation*}

We also have the following relations between the root vectors:\footnote{Similar relations can be found in the paper \cite{Xi1994}, but there appear to be some typos. They can also be derived from \eqref{RRel} by applying $\id \otimes B$ to both sides, for suitable linear maps $B$ on $\mathrm{End}\big(P_2V^{\otimes 2}\big)$. One can also check that the relations hold in the explicit representations in Remark~\ref{Remark1} below.}
For $i<l$ and $j<k$,
\begin{gather}\label{comm}
{E}_{il}{E}_{jk} =
\begin{cases}
{E}_{jk}{E}_{il}, & i< j < k < l \text{ or } i<l<j<k,\\
q^{-1} {E}_{jk}{E}_{il}, & i=j<k<l, \\
q {E}_{jk}{E}_{il}, & i<j<k=l,\\
{E}_{jk}{E}_{il} + \big(q-q^{-1}\big) {E}_{jl}{E}_{ik}, & i<j<l<k,
\end{cases}
\\
\label{comm1}
{E}_{li}{E}_{kj} =
\begin{cases}
{E}_{kj}{E}_{li}, & i< j < k < l \text{ or } i<l<j<k,\\
q^{} {E}_{kj}{E}_{li}, & i=j<k<l, \\
q^{-1} {E}_{kj}{E}_{li}, & i<j<k=l,\\
{E}_{kj}{E}_{li} - \big(q-q^{-1}\big) {E}_{lj}{E}_{ki}, & i<j<l<k.
\end{cases}
\end{gather}
By the first and fourth lines above,
\begin{gather}\label{comm2}
\hat{E}_{il}\hat{E}_{jk} + q^{-1} \hat{E}_{ik}\hat{E}_{jl} = \hat{E}_{jk}\hat{E}_{il} + q \hat{E}_{jl}\hat{E}_{ik}, \\
\hat{E}_{li}\hat{E}_{kj} + q^{} \hat{E}_{ki}\hat{E}_{lj} = \hat{E}_{kj}\hat{E}_{li} + q^{-1} \hat{E}_{lj}\hat{E}_{ki}. \label{comm3}
\end{gather}

For any $ \ii $ and $ \jj $ in $ \mathbb{N}^m $, define
\[
\hat{E}^{\pm}_{\ii \jj} :=
\begin{cases}
\hat{E}_{i_1j_1}\cdots \hat{E}_{i_mj_m}, & \pm(i_1-j_1),\dots,\pm(i_m-j_m)\geq 0,\\
0, & \text{else}.
\end{cases}
\]
 and let $ e_{ \jj \ii }:= e_{j_1i_1} \otimes \dots \otimes e_{j_mi_m}$. For $\ii, \jj \in \mathcal{W}_m$, define
\[
\tilde{E}^{\pm}_{\ii\jj} =
\sum_{\zeta \in D_{\ii}}q^{\pm(\inv(\tau) -\inv(\zeta))}\hat{E}^{\pm}_{\bar{\zeta}(\ii)\bar{\tau}(\jj)},
\]
where $\tau$ is an arbitrary element of $D_{\jj}$. Note that the notation does not depend on $\tau$, which is justified by the following lemma and the relation $\inv(\bar{\tau}) = (m-1)m/2 - \inv(\tau)$.

\begin{Lemma} \label{prop 1}\quad
	\begin{enumerate}\itemsep=0pt
		\item[$1.$] For every $\tau,\tau' \in D_A$, the following identity holds: \begin{equation*}%\label{key}
		\sum_{\sigma\in D_B} q^{\mp (\inv(\sigma)-\inv(\tau))} \hat{E}^{\pm}_{\tau(A) \sigma (B)} = \sum_{\sigma\in D_B} q^{\mp (\inv(\sigma)-\inv(\tau'))} \hat{E}^{\pm}_{\tau'(A) \sigma (B)} .
		\end{equation*} 	 	
		\item[$2.$] Likewise, for every $\sigma,\sigma' \in D_B$, \begin{equation*}%\label{key}
		\sum_{\tau\in D_A} q^{\mp (\inv(\tau)-\inv(\sigma))} \hat{E}^{\pm}_{\tau(A) \sigma (B)} = \sum_{\tau\in D_A} q^{\mp (\inv(\tau)-\inv(\sigma'))} \hat{E}^{\pm}_{\tau(A) \sigma' (B)} .
		\end{equation*} 	 	
	\end{enumerate}
	\end{Lemma}
\begin{proof}
	For $m=2$, both cases are equivalent to the relations in \eqref{comm} and \eqref{comm1}.
	
	Now suppose that $ m>2$. We only prove part 1, as part 2 is similar. By Lemma \ref{Gene}, it suffices to consider the case when $\tau'=s\tau$ where $s$ is a transposition, and assume without loss of generality that $\inv(\tau')= \inv(\tau)+1$. Define the two sets $D_B^{(1)}$ and $D_B^{(2)}$ by
	\begin{gather*}
		D_B^{(1)}= \{ \sigma \in D_B \colon s \sigma(B)=B \},\qquad
		D_B^{(2)}= \{ \sigma \in D_B \colon s \sigma(B) \neq B \}.
	\end{gather*}	
	First, note that by the second and third lines of \eqref{comm},
\[
	\sum_{\sigma\in D_B^{(1)}} q^{\inv(\sigma)-\inv(\tau)} \hat{E}^-_{\tau(A) \sigma (B)} = \sum_{\sigma\in D_B^{(1)}} q^{\inv(\sigma)-\inv(s\tau)} \hat{E}^-_{s \tau(A) \sigma (B)} .
\]
	Partition the set $D_B^{(2)}$ into two sets $D_-$ and $D_+$ of equal cardinality, where $\sigma \in D_-$ if and only if $\inv(\sigma) < \inv(s \sigma)$. Then
	\begin{gather*}
	\sum_{\sigma\in D_B^{(2)}} q^{\inv(\sigma)-\inv(\tau)} \hat{E}^-_{\tau(A) \sigma (B)} = \sum_{\sigma \in D_-} \big( q^{\inv(\sigma)-\inv(\tau)} \hat{E}^-_{\tau(A) \sigma (B)} + q^{\inv(\sigma)+1-\inv(\tau)} \hat{E}^-_{\tau(A) s\sigma (B)} \big) \\
\qquad \stackrel{\eqref{comm2}}{=} \sum_{\sigma \in D_-} \big( q^{\inv(\sigma)+1-\inv(\tau')} \hat{E}^-_{s\tau(A) s\sigma (B)} + q^{\inv(\sigma)-1-\inv(\tau)} \hat{E}^-_{s\tau(A) \sigma (B)} \big)\\
\qquad =\sum_{\sigma \in D_-} \big( q^{\inv(\sigma)+1-\inv(\tau')} \hat{E}^-_{\tau'(A) s\sigma (B)} + q^{\inv(\sigma)-\inv(\tau')} \hat{E}^-_{\tau'(A) \sigma (B)}\big)\\
\qquad =\sum_{\sigma\in D_B^{(2)}} q^{\inv(\sigma)-\inv(\tau')} \hat{E}^-_{\tau'(A) \sigma (B)} ,
	\end{gather*}
	as needed.
	The proof for $E^+$ is similar, where one uses \eqref{comm3} instead of \eqref{comm2}.
\end{proof}
